% =============================================================================
% exercises/ejercicios_cap02.tex
% Ejercicios propuestos — Capitulo 2: Procesos Electroquímicos
% =============================================================================

\section*{Ejercicios Propuestos}
\addcontentsline{toc}{section}{Ejercicios propuestos}
\label{sec:ejercicios_cap02}

\begin{bloque_ejercicios}[Ejercicios -- Capitulo 2: Procesos Electroqu\'{\i}micos]

\begin{ejercicios}

% ----- Ejercicio 1: Números de oxidación ------------------------------------
\item Determine el n\'umero de oxidaci\'on del elemento subrayado en cada
uno de los siguientes compuestos o iones, aplicando las reglas del
\textbf{Cuadro~\ref{tab:5}}:

{%
\centering
\begin{tabular}{llll}
\toprule
\textbf{Especie} & \textbf{Elemento} & \textbf{Especie} & \textbf{Elemento} \\
\midrule
\ce{H2\underline{S}O4}   & S & \ce{\underline{Mn}O4^-}      & Mn \\
\ce{\underline{Cr}2O7^{2-}} & Cr & \ce{\underline{N}O3^-}    & N  \\
\ce{\underline{P}O4^{3-}}   & P  & \ce{\underline{Cl}O3^-}   & Cl \\
\ce{K2\underline{Cr}O4}   & Cr & \ce{Na2\underline{S}2O3}    & S  \\
\bottomrule
\end{tabular}
\par
}%

\textit{[Respuesta: S en \ce{H2SO4}, $+6$;
                    Mn en \ce{MnO4^-}, $+7$;
                    Cr en \ce{Cr2O7^{2-}}, $+6$;
                    N en \ce{NO3^-}, $+5$;
                    P en \ce{PO4^{3-}}, $+5$;
                    Cl en \ce{ClO3^-}, $+5$;
                    Cr en \ce{K2CrO4}, $+6$;
                    S en \ce{Na2S2O3}, $+2$]}

% ----- Ejercicio 2: Identificar agente oxidante y reductor ------------------
\item En cada una de las siguientes reacciones redox, identifique el
\textbf{agente oxidante} y el \textbf{agente reductor}, e indique qu\'e
especie se oxida y cu\'al se reduce:

\begin{enumerate}[label=\alph*)]
    \item \ce{Zn (s) + 2HCl (ac) -> ZnCl2 (ac) + H2 (g)}
    \item \ce{2Fe^2+ (ac) + Cl2 (g) -> 2Fe^3+ (ac) + 2Cl^- (ac)}
    \item \ce{MnO4^- + 5Fe^2+ + 8H+ -> Mn^2+ + 5Fe^3+ + 4H2O}
    \item \ce{2Al (s) + Fe2O3 (s) -> Al2O3 (s) + 2Fe (s)}
\end{enumerate}

% ----- Ejercicio 3: Potencial de celda y espontaneidad ----------------------
\item Usando los potenciales de reducci\'on est\'andar del
\textbf{Cuadro~\ref{tabla7}}, calcule el potencial de celda $E^\circ_{\text{celda}}$
para cada sistema y determine si la reacci\'on es espont\'anea:

\begin{enumerate}[label=\alph*)]
    \item \ce{Zn (s) | Zn^2+ (ac) || Ag+ (ac) | Ag (s)}
    \item \ce{Cu (s) | Cu^2+ (ac) || Fe^3+ (ac) | Fe (s)}
    \item \ce{Mg (s) | Mg^2+ (ac) || Al^3+ (ac) | Al (s)}
    \item \ce{Fe (s) | Fe^2+ (ac) || Cu^2+ (ac) | Cu (s)}
\end{enumerate}

\textit{[Respuesta:
(a) $E^\circ = 0.80 - (-0.76) = \SI{1.56}{\volt}$, espont\'anea;
(b) $E^\circ = -0.04 - 0.34 = \SI{-0.38}{\volt}$, no espont\'anea;
(c) $E^\circ = -1.76 - (-2.36) = \SI{0.60}{\volt}$, espont\'anea;
(d) $E^\circ = 0.34 - (-0.44) = \SI{0.78}{\volt}$, espont\'anea]}

% ----- Ejercicio 4: Diagrama de celda ---------------------------------------
\item Para la reacci\'on global:
\[
    \ce{2Al (s) + 3Cu^2+ (ac) -> 2Al^3+ (ac) + 3Cu (s)}
\]

\begin{enumerate}[label=\alph*)]
    \item Escriba las semirreacciones de oxidaci\'on y de reducci\'on
          por separado.
    \item Dibuje el diagrama de celda usando la notaci\'on convencional
          (\'anodo a la izquierda, c\'atodo a la derecha, $|$ para fronteras
          de fase y $||$ para el puente salino).
    \item Calcule $E^\circ_{\text{celda}}$ e indique si la reacci\'on
          es espont\'anea.
    \item ?`Cu\'al electrodo es el \'anodo y cu\'al es el c\'atodo?
          ?`En qu\'e direcci\'on fluyen los electrones en el circuito externo?
\end{enumerate}

\textit{[Respuesta:
(b) \ce{Al(s)|Al^3+(ac)||Cu^2+(ac)|Cu(s)};
(c) $E^\circ = 0.34 - (-1.76) = \SI{2.10}{\volt}$, espont\'anea;
(d) Al es el \'anodo, Cu es el c\'atodo; los electrones fluyen de Al hacia Cu.]}

% ----- Ejercicio 5: Relación con energía libre de Gibbs ---------------------
\item La relaci\'on entre el potencial de celda y la energ\'{\i}a libre de
Gibbs es:
\[
    \Delta G^\circ = -nFE^\circ_{\text{celda}}
\]
donde $n$ es el n\'umero de moles de electrones transferidos y
$F = \SI{96485}{\coulomb\per\mole}$ es la constante de Faraday.

Para la celda de Daniell \ce{Zn(s)|Zn^2+(ac)||Cu^2+(ac)|Cu(s)}
con $E^\circ_{\text{celda}} = \SI{1.10}{\volt}$:

\begin{enumerate}[label=\alph*)]
    \item Determine el n\'umero de electrones transferidos $n$ a partir
          de las semirreacciones balanceadas.
    \item Calcule $\Delta G^\circ$ de la reacci\'on en \si{\kilo\joule\per\mole}.
    \item ?`Es consistente el signo de $\Delta G^\circ$ con el valor
          positivo de $E^\circ_{\text{celda}}$? Explique.
    \item Calcule la constante de equilibrio $K$ a \SI{298}{\kelvin}
          usando $\Delta G^\circ = -RT\ln K$.
\end{enumerate}

\textit{[Respuesta:
(a) $n = 2$;
(b) $\Delta G^\circ = -(2)(96485)(1.10) = \SI{-212.3}{\kilo\joule\per\mole}$;
(c) S\'{\i}, $\Delta G^\circ < 0$ indica proceso espont\'aneo, consistente con
    $E^\circ > 0$;
(d) $\ln K = 212300/(8.314 \times 298) \Rightarrow K \approx 1.5 \times 10^{37}$]}

% ----- Ejercicio 6: Corrosión y prevención ----------------------------------
\item Con base en los potenciales de reducci\'on del
\textbf{Cuadro~\ref{tabla7}} y los conceptos de corrosi\'on estudiados:

\begin{enumerate}[label=\alph*)]
    \item El acero (hierro) de un barco se protege uniendo placas de zinc
          a su casco. Explique por qu\'e el zinc funciona como \textbf{\'anodo
          de sacrificio} y el hierro queda protegido.  Apoye su respuesta
          con los potenciales de reducci\'on de ambos metales.
    \item ?`Ser\'{\i}a m\'as adecuado usar cobre en lugar de zinc como \'anodo
          de sacrificio para proteger el hierro? Justifique su respuesta.
    \item Una tuber\'{\i}a de hierro est\'a conectada accidentalmente a una
          de cobre. Calcule $E^\circ_{\text{celda}}$ del par galv\'anico que
          se forma e identifique cu\'al metal se corroe.
    \item Mencione y describa brevemente \textbf{dos m\'etodos} adicionales
          (distintos al \'anodo de sacrificio) para prevenir la corrosi\'on
          del hierro.
\end{enumerate}

\textit{[Respuesta:
(c) $E^\circ = 0.34 - (-0.44) = \SI{0.78}{\volt}$; el hierro es el \'anodo
    y se corroe porque tiene menor potencial de reducci\'on que el cobre.]}

% ----- Ejercicio 7 (avanzado): Celda con tres electrodos --------------------
\item \textbf{[Avanzado]} Se construye una celda con los siguientes tres
pares redox: \ce{MnO4^-}/\ce{Mn^2+} ($E^\circ = \SI{1.51}{\volt}$),
\ce{Fe^3+}/\ce{Fe^2+} ($E^\circ = \SI{0.77}{\volt}$) y
\ce{Zn^2+}/\ce{Zn} ($E^\circ = \SI{-0.76}{\volt}$).

\begin{enumerate}[label=\alph*)]
    \item Ordene los tres pares de mayor a menor potencial de reducci\'on
          e identifique cu\'al act\'ua como c\'atodo y cu\'al como \'anodo
          en la celda m\'as energ\'etica posible.
    \item Escriba la reacci\'on global balanceada para la celda de mayor
          $E^\circ_{\text{celda}}$.
    \item Calcule $\Delta G^\circ$ para esta celda, sabiendo que se
          transfieren $n = 10$ moles de electrones.
    \item ?`Qu\'e papel juega el par \ce{Fe^3+}/\ce{Fe^2+} si se intercala
          entre los otros dos electrodos?
\end{enumerate}

\end{ejercicios}

\end{bloque_ejercicios}
% =============================================================================
% FIN DE ejercicios_cap02.tex
% =============================================================================
